\section{Introduction}

\subsection{Problem Statement}

\subsubsection{Software Testing Perspective}

When testing code consuming numerical data, it is common to use random transforms to produce synthetic data sets for testing.

However, this introduces the probability that any given run of the test will fail to catch edge cases, or fail due to a specific random seed.

The problem we address in this paper is: \textit{How can we efficiently and reliably estimate the required number of test runs needed to establish the absence of failures with a certain confidence level, given an unknown relationship between the input data and test failure, and knowledge of the distribution of the input data?}

\subsubsection{Mathematical Setup}

    \subsubsubsection{Population and Sets}

    Each input to the function under test consists of a fixed-size set of elements $X^{(j)}$.

    Let $P = \{ X^{(1)}, X^{(2)}, \dots \}$ be a countably infinite population of sets, representing all possible inputs for the function under test $F$. 
    
    \[
        X^{(j)} = \{ X_0^{(j)}, X_1^{(j)}, \dots, X_{n-1}^{(j)} \}, \quad j \in \mathbb{N}.
    \]

    Each element $X_i^{(j)}$ is drawn i.i.d. and has probability $p$ of being an \textbf{outlier} (an event we denote $O_i^{(j)} = 1$ if $X_i^{(j)}$ is an outlier, $0$ otherwise):
    \[
        \Pr[ O_i^{(j)} = 1 ] = p
    \]

    Where an \textbf{outlier} is defined as an element exhibiting characteristics associated with an increased likelihood of test failure.

    \subsubsubsection{Blackbox Function}

    There is a function under test, which we treat as a blackbox function
    \[
        F : \mathcal{X} \to \{0,1\}, \quad F(X^{(j)}) =
        \begin{cases}
        1 & \text{if set passes} \\
        0 & \text{if set fails}
        \end{cases}
    \]

    We assume there exists a relationship between the presence of outliers in ( $X^{(j)}$ ) and failure, but the exact nature of this relationship is unknown:
    \[
        \text{unknown mapping } F(X^{(j)}) \sim f(O_0^{(j)}, \dots, O_{n-1}^{(j)})
    \]

    \subsubsubsection{Sampling Problem}

    Let ( S \subseteq P ) be the subset of sets actually observed/tested.

    Define the event
    \[
        E = \exists X^{(j)} \in S : F(X^{(j)}) = 0
    \]

    i.e., "we observe at least one failing set."

    The goal is:

    \[
        \text{ Find the minimal } |S| \text{ such that } \Pr[E] \ge \gamma
    \]

    where ( \gamma \in (0,1) ) is a target confidence (e.g., 0.95).

    Corresponding to the problem of finding the minimum number of samples needed to have a confidence of at least ( \gamma ) that we have observed at least one failing set, given the unknown relationship between outliers and failure, and the probability of outliers in the population.

\subsection{Related Work}

Common approaches to mitigate this risk include fixed-budget testing, empirical convergence, fuzzing, or Monte Carlo estimation, all of which have their own drawbacks and limitations which limit their applicability to this problem. One major limitation of many approaches is the compute or time cost of running the test, making their use in CI pipelines or other time-sensitive contexts difficult

\subsection{Comparison to Classical Testing Methods}

\subsubsubsection{Fixed-budget testing}

\begin{itemize}
    \item \textit{Approach} - Select a large fixed number of iterations $k$ based on heuristics, and run the test for that many iterations \cite{myers2011art}.
    \item \textit{Advantages} - This method is simple and easy to implement
    \item \textit{Limitations} - It does not provide any guarantees about the confidence level of detecting failure, and can be computationally expensive if $k$ is set too high. Additionally, in the case of the absence of failure, this method provides no ability to stop early, and thus can waste resources.
\end{itemize}

\subsubsubsection{Empirical convergence / burn-in heuristics}

\begin{itemize}
    \item \textit{Approach} - Test until the observed pass ratio converges to a stable value, or until a certain number of consecutive passes are observed \cite{myers2011art} \footnote{Myers et al. describe stopping rules based on reaching estimated error thresholds or elapsed testing time, which is conceptually analogous to monitoring convergence of observed pass/fail ratios.}.
    \item \textit{Advantages} - This method can be more efficient than fixed-budget testing, as it allows for early stopping if the test appears to be consistently passing.
    \item \textit{Limitations} - It can be sensitive to the choice of convergence criteria, and may not provide guarantees about the confidence level of detecting failure. Additionally, very rare failures may not be observed within a reasonable number of iterations, leading to false confidence in the reliability of the code.
\end{itemize}

\subsubsubsection{Worst-case / adversarial thinking}

\begin{itemize}
    \item \textit{Approach} - Design tests to specifically target edge cases and potential failure modes, rather than relying on random sampling.
    \item \textit{Advantages} - This can be effective in catching specific known issues, and can be more efficient than random sampling if the failure modes are well understood.
    \item \textit{Limitations} - It may not provide a comprehensive assessment of the code's reliability, and can be time-consuming to design and implement. Additionally, it may not be effective in catching unknown failure modes or edge cases that were not anticipated during test design.
\end{itemize}

\subsubsubsection{Classical statistical power analysis}

\begin{itemize}
    \item \textit{Approach} - Calculate the required sample size to achieve a desired power level for detecting a specific effect size, based on assumptions about the underlying distribution of the data and the expected failure rate \cite{cohen1988statistical}.
    \item \textit{Advantages} - This can provide a more rigorous basis for determining the number of iterations required, and can be tailored to specific tests and failure modes.
    \item \textit{Limitations} - It relies heavily on the accuracy of the assumptions made, and may not be robust to mis-specification of parameters. Additionally, it may not be practical in cases where the underlying distribution is unknown or difficult to model and is specific to the particular test being run, making it less generalizable.
\end{itemize}

\subsubsection{Comparison to Recent Methods}

\subsubsubsection{Thing A}
blah blah

\subsubsubsection{Thing B}
blah blah